\documentclass[11pt]{article}

\usepackage[margin=1in]{geometry}
\usepackage{amsmath,amssymb,amsthm}
\usepackage{mathtools}
\usepackage{enumitem}
\usepackage[colorlinks=true,linkcolor=blue,citecolor=blue,urlcolor=blue]{hyperref}

\newtheorem{theorem}{Theorem}
\newtheorem{lemma}[theorem]{Lemma}
\newtheorem{proposition}[theorem]{Proposition}
\newtheorem{corollary}[theorem]{Corollary}
\theoremstyle{definition}
\newtheorem{remark}[theorem]{Remark}

\DeclareMathOperator{\lcm}{lcm}
\newcommand{\Z}{\mathbb{Z}}
\newcommand{\Q}{\mathbb{Q}}
\newcommand{\R}{\mathbb{R}}
\newcommand{\N}{\mathbb{N}}
\newcommand{\dn}{d_n}
\newcommand{\code}[1]{\texttt{#1}}
\newcommand{\Rn}{R_n}

\title{A formal proof, in Lean~4, that one of\\
  $\zeta(5),\zeta(7),\dots,\zeta(33)$ is irrational}

\author{Formalization of W.~Zudilin's elementary argument\\
  \small(arXiv:1801.09895), carried out in Lean~4 / Mathlib}

\date{July 2026}

\begin{document}
\maketitle

\begin{abstract}
We report a complete, machine-checked formalization, in the Lean~4 proof
assistant on top of Mathlib, of a theorem of the shape ``at least one of the
odd zeta values $\zeta(5),\zeta(7),\dots,\zeta(s)$ is irrational.'' The
argument formalized is the elementary one of W.~Zudilin
\cite{Zudilin2018}, whose only analytic inputs are Stirling's formula and a
bound on the least common multiple $d_n=\lcm(1,\dots,n)$. Because the current
Mathlib does not contain the prime number theorem, we cannot use Zudilin's
headline value $s=25$ (which needs $d_n^{1/n}\to e$); instead we formalize,
from scratch and sorry-free, Hanson's elementary bound $d_n\le 3^n$ and take
$s=33$ --- precisely the fallback that Zudilin himself records. The final
theorem
\[
  \exists\, j\in\{5,7,9,\dots,33\},\qquad \zeta(j)\notin\Q,
\]
is proved with no axioms beyond Lean's three standard ones
(\code{propext}, \code{Classical.choice}, \code{Quot.sound}); in particular it
uses no \code{sorry} and no kernel-external computation
(\code{native\_decide}). We describe the construction, the four analytic and
arithmetic ingredients, the reconstruction of Hanson's bound, and the
engineering of the development.
\end{abstract}

\tableofcontents

\section{Introduction}

The irrationality of $\zeta(s)=\sum_{n\ge1}n^{-s}$ is classical for even $s$
and, after Ap\'ery, for $s=3$; for odd $s\ge5$ only partial results are known.
The prototypical such result --- ``infinitely many odd zeta values are
irrational,'' and its finite refinements ``at least one of
$\zeta(5),\dots,\zeta(s)$ is irrational'' --- is due to Ball--Rivoal
\cite{BallRivoal} and Rivoal, and traditionally rests on the saddle-point
method and/or linear-independence criteria, both regarded as decidedly
non-elementary. Zudilin \cite{Zudilin2018} gave, ``by elementary means,'' the
first proof of such a statement whose only technical ingredients are the prime
number theorem and Stirling's approximation. The present work is a full formal
verification of that argument.

\subsection*{The formalized statement}

Let, in Lean,
\[
  \code{zetaVal}\ j \;=\; \sum_{n=0}^{\infty}\frac{1}{(n+1)^{j}}
  \;=\;\zeta(j),\qquad
  \code{oddIdx}\;=\;\{\,j\in\N : 5\le j\le 33,\ j\ \text{odd}\,\}.
\]
Our main theorem is
\begin{equation}\label{eq:main}
  \code{zeta\_odd\_irrational} \;:\;
  \exists\, j\in\code{oddIdx},\ \operatorname{Irrational}(\code{zetaVal}\,j),
\end{equation}
i.e.\ \emph{at least one of $\zeta(5),\zeta(7),\dots,\zeta(33)$ is
irrational}. Here $\operatorname{Irrational}$ is Mathlib's predicate
$x\notin\operatorname{range}(\Q\hookrightarrow\R)$. A \code{\#print axioms}
query returns
\[
  \text{\code{[propext, Classical.choice, Quot.sound]}}
\]
and the whole development compiles with no \code{sorry}.

\subsection*{Why $s=33$ and not $s=25$}

Zudilin's optimal elementary result is for $s=25$. Its proof uses the prime
number theorem in the equivalent form $\lim_{n\to\infty}d_n^{1/n}=e$. Mathlib
(at the commit used here, on \code{leanprover/lean4:v4.33.0-rc1}) does
\emph{not} contain the prime number theorem; its sharpest built-in bound is the
Chebyshev estimate $\psi(x)\le x\log 4+o(x)$, giving only
$\limsup_n d_n^{1/n}\le 4$, which is too weak for \emph{every} $s$. The
external \emph{PrimeNumberTheoremAnd} project would pin an incompatible Mathlib
revision and cannot be added without a forbidden dependency update. We
therefore follow Zudilin's own remark (\cite{Zudilin2018}, \S4): using the
weaker, \emph{elementary} bound $d_n<3^n$ of Hanson \cite{Hanson1972} and the
choice $s=33$ ``to arrive at the same conclusion.'' Reconstructing Hanson's
bound in Lean turned out to be one of the substantial pieces of the project
(\S\ref{sec:hanson}).

\section{The construction}\label{sec:construction}

Fix an odd $s\ge7$ and write $q=(s+1)/2\ge4$; the formalization takes $s=33$,
$q=17$, but all of the asymptotic analysis is developed for general $q$. Set
\begin{equation}\label{eq:R}
  \Rn(t)=
  \frac{n!^{\,s-5}\,2^{6n}\prod_{j=1}^{n}(t-j)\prod_{j=1}^{n}(t+n+j)
        \prod_{j=1}^{3n}\!\bigl(t-n-\tfrac12+j\bigr)}
       {\prod_{j=0}^{n}(t+j)^{s}} .
\end{equation}
The two families of linear forms are
\[
  r_n=\sum_{\nu=1}^{\infty}\Rn(\nu),\qquad
  \hat r_n=\sum_{\nu=1}^{\infty}\Rn\!\bigl(\nu-\tfrac12\bigr).
\]
Since $\Rn$ vanishes at $1,\dots,n$ and at $\tfrac12,\dots,n-\tfrac12$, both
sums may be started at $\nu=n+1$; the corresponding summands, written in
factorial form, are the Lean definitions \code{c q n k}$=\Rn(n+1+k)$ and
\code{chat q n k}$=\Rn(n+\tfrac12+k)$, and
\[
  \code{r q n}=\sum_{k\ge0}\code{c q n k},\qquad
  \code{rhat q n}=\sum_{k\ge0}\code{chat q n k}.
\]

The proof has four ingredients, mirrored by the module layout:
\begin{enumerate}[label=(\Alph*)]
\item \emph{Arithmetic} (\S\ref{sec:arith}, \code{Forms.lean}): $r_n$ and
  $\hat r_n$ are $\Z[1/d_n]$-linear combinations of $1,\zeta(3),\zeta(5),\dots,
  \zeta(s)$, and the ``twist by half'' makes their $\zeta(3)$-coefficients
  proportional, so that $7r_n-\hat r_n$ lives in
  $\Q+\Q\zeta(5)+\dots+\Q\zeta(s)$.
\item \emph{Asymptotics} (\S\ref{sec:asymp}, Lemma~4): $r_n^{1/n},\hat
  r_n^{1/n}\to g(x_0)$ and $r_n/\hat r_n\to1$.
\item \emph{The denominator bound} (\S\ref{sec:hanson}, \code{DnBound.lean}):
  $d_n\le 3^n$.
\item \emph{The numerical estimate} (\S\ref{sec:numeric},
  \code{Numeric.lean}): $3^{33}\,g(x_0)<1$.
\end{enumerate}
Assembling these yields the contradiction (\S\ref{sec:main}).

\section{Arithmetic of the linear forms}\label{sec:arith}

This is the combinatorial heart of the development
(\code{Zeta5Odd/Forms.lean}). We reproduce Zudilin's Lemmas 1--3 and the
$\zeta(3)$-elimination.

\subsection{Partial fractions and integrality (Lemmas 1--2)}

Write $\Rn(t)=\sum_{i=1}^{s}\sum_{k=0}^{n}a_{i,k}/(t+k)^i$ (its partial-fraction
decomposition, eq.~(e04) of \cite{Zudilin2018}). Two facts are needed.

\begin{lemma}[Integrality; Lemma~1 of \cite{Zudilin2018}]\label{lem:int}
$d_n^{\,s-i}\,a_{i,k}\in\Z$ for $1\le i\le s$, $0\le k\le n$.
\end{lemma}

\begin{lemma}[Well-poised symmetry; Lemma~2]\label{lem:sym}
$a_{i,k}=(-1)^{i-1}a_{i,n-k}$; hence $\sum_k a_{i,k}=0$ for even $i$.
\end{lemma}

In Lean these are packaged as
\[
  \code{partialFraction\_exists}\ n :\ \exists\, a:\N\to\N\to\R,\
  (\text{decomposition})\ \wedge\ (\text{integrality})\ \wedge\ (\text{symmetry}).
\]
The three conjuncts are proved as follows.

\emph{Decomposition and integrality} (\code{pf\_decomp}) follow Zudilin's route
of writing $\Rn$ as a product of the six ``simpler'' rational functions of
\cite{Zudilin2018} (eqs.\ before (e04)), each of which has denominator
$\prod_{j=0}^n(t+j)$ and \emph{integer} partial-fraction coefficients given by
explicit binomials, e.g.
\[
  \frac{n!}{\prod_{j=0}^n(t+j)}=\sum_{k=0}^n\frac{(-1)^k\binom nk}{t+k},
  \quad
  \frac{2^{2n}\prod_{j=1}^n(t-\tfrac12+j)}{\prod_{j=0}^n(t+j)}
   =\sum_{k=0}^n\frac{\binom{2k}{k}\binom{2n-2k}{n-k}}{t+k}.
\]
Each of these six identities is proved (lemmas \code{f1\_id},\dots,\code{f6\_id})
via a reusable ``clear the denominator and interpolate'' engine
\code{pf\_general}, which reduces the claim to a polynomial identity of degree
$<n+1$ holding at the $n+1$ nodes $t=-k$ (Mathlib's
\code{Polynomial.eq\_of\_natDegree\_lt\_card\_of\_eval\_eq}). The product of
$28$ copies of the first simple function with the remaining five reproduces
$\Rn$ (lemma \code{Rn\_as\_simpleProd}, matching the $n!^{\,28}$, $2^{6n}$ and
denominator power $s=33=28+5$). Integrality of the product's coefficients is an
induction on the number of factors (\code{pf\_prod}, step \code{pf\_mul\_simple})
whose analytic engine is the two-pole identity
\[
  \frac{1}{u^{\,i}(u+\delta)}
    =\sum_{r<i}\frac{(-1)^r}{\delta^{\,r+1}\,u^{\,i-r}}
     +\frac{(-1)^i}{\delta^{\,i}(u+\delta)}
  \qquad(\code{pf\_two\_pole}),
\]
applied with $u=t+k$, $\delta=(k'-k)$; the $\delta^{-p}$ factors are cleared by
$d_n^{\,p}$ because $|k'-k|\le n$ divides $d_n$
(\code{dvd\_lcmUpto}). Uniqueness of the coefficients (\code{pf\_unique}),
needed to transport the symmetry, is proved by a residue-peeling limit
argument: near a pole $-k_0$ the scaled difference
$(t+k_0)^{P}\sum(a-b)_{i,k}/(t+k)^i$ tends both to $0$ and to
$(a-b)_{P,k_0}$.

\emph{Symmetry} (Lemma~\ref{lem:sym}) is derived from the well-poised
functional equation
\[
  \Rn(-t-n)=-\Rn(t)\qquad(\code{Rn\_wellPoised}),
\]
proved directly from \eqref{eq:R} by reindexing the four finite products (the
global sign $-1$ coming from $s$ being odd), combined with uniqueness.

\subsection{The zeta representation (Lemma~3)}

\begin{lemma}[Lemma~3 of \cite{Zudilin2018}]\label{lem:l3}
With $a_i=\sum_k a_{i,k}$,
\[
  r_n=\sum_{\substack{i=3\\ i\ \mathrm{odd}}}^{s} a_i\,\zeta(i)+a_0,
  \qquad
  \hat r_n=\sum_{\substack{i=3\\ i\ \mathrm{odd}}}^{s}
     a_i\,(2^{i}-1)\,\zeta(i)+\hat a_0,
\]
with $d_n^{\,s-i}a_i\in\Z$ and $d_n^{\,s}a_0,\ d_n^{\,s}\hat a_0\in\Z$.
\end{lemma}

The Lean statement \code{repr\_combined} multiplies through by $d_n^{s}$ and
plumbs the integer constants through existentials. The $r_n$ half (eq.~(e07))
uses $\sum_{\nu}(\nu+k)^{-i}=\zeta(i)-\sum_{\ell\le k}\ell^{-i}$
(\code{tsum\_shift\_zeta}); a $\code{tsum}$/finite-sum interchange after
discarding the even columns (Lemma~\ref{lem:sym}) and the $i=1$ column
($\sum_k a_{1,k}=0$, from convergence of the residue sum); and the harmonic
integrality $d_n^{\,i}\sum_{\ell\le k}\ell^{-i}\in\Z$
(\code{harmonic\_integrality}), which is absorbed into $a_0$. The $\hat r_n$
half (eq.~(e08)) is the analytic subtlety of the paper: $\Rn$ does not vanish
at the half-integer nodes, so the summation is shifted to start at
$\nu=-m-\tfrac12$ with $m=\lfloor(n-1)/2\rfloor$ (\code{rhat\_shift}); the
resulting half-integer harmonic heads have denominators $\le n$ and are cleared
by $d_n$ and $d_{n-1}$ (\code{odd\_harmonic\_integrality},
\code{signed\_harmonic\_integrality}); the identity
$\sum_{\ell\ge1}(\ell-\tfrac12)^{-i}=(2^i-1)\zeta(i)$ supplies the coefficient
$(2^i-1)$.

\subsection{Elimination of $\zeta(3)$}

Because the $\zeta(i)$-coefficient of $\hat r_n$ is $a_i(2^i-1)$, that of
$7r_n-\hat r_n$ is $a_i\bigl(7-(2^i-1)\bigr)=a_i(8-2^i)$, which
\emph{vanishes at $i=3$} (as $2^3-1=7$). Hence, for $s=33$,
\begin{equation}\label{eq:elim}
  \code{elim\_integer}\ n:\quad
  d_n^{\,33}\,(7r_n-\hat r_n)
    =\sum_{j\in\code{oddIdx}} A_j\,\zeta(j)+A_0,\qquad A_j,A_0\in\Z,
\end{equation}
where $A_j=d_n^{\,j}\,a_j\,(8-2^{j})$ and $A_0=7a_0-\hat a_0$, both integral by
Lemma~\ref{lem:l3}. This is the exact interface consumed by the endgame.

\section{Asymptotics (Lemma~4)}\label{sec:asymp}

The asymptotic analysis (modules \code{Basic}, \code{Localize}, \code{Window},
\code{Roots}, \code{Ratio}) is developed for all $q\ge4$. Writing
\[
  f(x)=\frac{x+3}{x}\Bigl(\frac{x+1}{x+2}\Bigr)^{q},\qquad
  g(x)=\frac{2^6(x+3)^6(x+1)^{s+1}}{(x+2)^{2(s+1)}},
\]
and letting $x_0\in(0,1)$ be the unique positive solution of $f(x)=1$, one has
\[
  \code{tendsto\_root\_r},\ \code{tendsto\_root\_rhat}:\quad
  r_n^{1/n},\ \hat r_n^{1/n}\longrightarrow g(x_0),
  \qquad
  \code{tendsto\_ratio}:\quad r_n/\hat r_n\to1 .
\]
The proof localizes both series to a window $|k-x_0n|\le\varepsilon n$, controls
the termwise ratio $c_{k+1}/c_k\sim f(k/n)^2$ via a quantitative Stirling
estimate, and passes to $n$-th roots, all of which absorb the subexponential
prefactors. This part predates the present campaign and is used as a black box;
it supplies in particular
\[
  \code{seven\_r\_sub\_rhat\_pos\_eventually}:\quad
  \forall^{\infty}n,\ 7r_n-\hat r_n>0 .
\]

\section{The denominator bound $d_n\le 3^n$}\label{sec:hanson}

Since Mathlib lacks the prime number theorem, the bound $d_n\le3^n$
(module \code{Zeta5Odd/DnBound.lean}) had to be built from the ground up. This
is an elementary but genuinely intricate theorem of Hanson \cite{Hanson1972};
our formalization is, to our knowledge, new.

\subsection{Divisibility}

Let $a_0=2,\ a_1=3,\ a_2=7,\ a_3=43,\dots$ be the Sylvester sequence
($a_{k+1}=a_0a_1\cdots a_k+1$; \code{sylv}, \code{sylvProd}) and set
\[
  C_n=\frac{n!}{\prod_{i\ge0}\bigl(\lfloor n/a_i\rfloor!\bigr)}
  \qquad(\code{hansonC}).
\]
Only finitely many factors are non-trivial, since $a_i>n$ forces
$\lfloor n/a_i\rfloor=0$. One shows $C_n\in\Z$
(\code{hansonDenom\_dvd}), and, crucially,
\[
  d_n=\lcm(1,\dots,n)\ \big|\ C_n
  \qquad(\code{lcmUpto\_dvd\_hansonC}),
\]
by comparing $p$-adic valuations through Legendre's formula and the arithmetic
inequality $\sum_i\lfloor m/a_i\rfloor\le m-1$ (\code{core\_sum\_le}), itself a
consequence of $\sum_i a_i^{-1}=1-1/(a_0\cdots a_{N-1})<1$
(\code{sum\_inv\_sylv}).

\subsection{Size}

It remains to bound $C_n\le3^n$, equivalently
$\log n!\le n\log3+\sum_i\log\lfloor n/a_i\rfloor!$. The margin is thin: the
limiting exponential rate is
$\sum_i (\log a_i)/a_i=1.08239\ldots<\log3=1.09861\ldots$, a gap of only
$0.0162$ per unit $n$. The Lean proof splits at $N_0=1600$:
\begin{itemize}[nosep]
\item for $n<N_0$ the bound is reduced (via the truncation
  $C_n=n!/\prod_{i<4}\lfloor n/a_i\rfloor!$, valid for $n<a_4=1807$) to a
  finite $\N$-inequality and discharged by the Lean kernel's \code{decide}
  (GMP arithmetic; \emph{no} \code{native\_decide}, hence no extra axiom);
\item for $n\ge N_0$ we combine sharp two-sided Stirling bounds
  (\code{log\_factorial\_le/ge}), the tangent line of the convex $t\mapsto
  t\log t$, and a \emph{certified rational} bound
  $\sum_i(\log a_i)/a_i<\log3$ (obtained from integer inequalities such as
  $7^{26}<2^{73}$, $43^{7}<2^{38}$ and a geometric tail), through a
  self-bootstrapping induction on the Sylvester index (\code{hanson\_key}).
\end{itemize}
The whole file certifies
$\code{lcmUpto\_le\_three\_pow}:\ (d_n:\R)\le3^n$ with axioms
\code{[propext, Classical.choice, Quot.sound]}.

\section{The numerical estimate}\label{sec:numeric}

At $x_0$ one has $x_0=2.2893\ldots\times10^{-5}$ and $\log g(x_0)=-36.3834\ldots$,
while $3^{33}=e^{36.2543\ldots}$, so $3^{33}g(x_0)=e^{-0.129}<1$. The Lean
lemma \code{g\_small} proves $3^{33}\,g(17,x)<1$ for the (unique) positive $x$
with $f(17,x)=1$. Two ingredients: (i) $x\le\tfrac1{40}$, by a
quadratic--Bernoulli argument reducing $f(x)=1$ to a factored polynomial
inequality; and (ii) $g$ increasing near $0$ (via the sign of $(\log g)'$), so
$g(x)\le g(\tfrac1{40})$, together with the \emph{exact rational} inequality
\[
  3^{33}g\bigl(17,\tfrac1{40}\bigr)<1
  \iff 64\cdot121^{6}\cdot41^{34}\cdot40^{28}<3^{239},
\]
closed by \code{norm\_num}.

\section{Assembly and the contradiction}\label{sec:main}

The endgame (\code{Zeta5Odd/Main.lean}) is Zudilin's final page. Suppose, for
contradiction, that every $\zeta(j)$ with $j\in\code{oddIdx}$ is rational; a
common denominator $a\in\N_{>0}$ then satisfies $a\,\zeta(j)\in\Z$
(\code{exists\_common\_denom}). Put
$b_n=d_n^{\,33}(7r_n-\hat r_n)$. From \eqref{eq:elim},
\[
  a\,b_n=\sum_{j\in\code{oddIdx}}A_j\,(a\zeta(j))+a\,A_0\in\Z .
\]
On the other hand $b_n\ge0$ and, by \S\ref{sec:asymp}--\ref{sec:numeric}, the
$n$-th root of $b_n$ satisfies
\[
  b_n^{1/n}
    =\bigl(d_n^{1/n}\bigr)^{33}(7r_n-\hat r_n)^{1/n}
    \le 3^{33}\,(7r_n-\hat r_n)^{1/n}
    \longrightarrow 3^{33}g(x_0)<1,
\]
using $d_n\le3^n$ and $(7r_n-\hat r_n)^{1/n}\to g(x_0)$
(\code{tendsto\_seven\_root}, deduced from Lemma~4). Hence a root test gives
$b_n\to0$, so $a\,b_n\to0$; but $7r_n-\hat r_n>0$ eventually makes $a\,b_n$ a
\emph{positive integer}, so $a\,b_n\ge1$ --- a contradiction. This is
\eqref{eq:main}.

\section{Verification and engineering}

\paragraph{Toolchain.} Lean~\code{v4.33.0-rc1}, Mathlib at commit
\code{cd580e54} (master). The development is one Lake library, \code{Zeta5Odd},
of ten modules.

\paragraph{Trust.} The final theorem \code{zeta\_odd\_irrational} depends only
on \code{[propext, Classical.choice, Quot.sound]}. There is no \code{sorry}
anywhere in the library, and no use of \code{native\_decide} (which would add
\code{Lean.ofReduceBool}); the one finite case analysis (\S\ref{sec:hanson})
uses kernel \code{decide}.

\paragraph{Module map.}
\begin{center}
\begin{tabular}{ll}
\code{Basic},\,\code{Localize},\,\code{Window},\,\code{Roots},\,\code{Ratio}
  & Lemma~4 asymptotics.\\
\code{ZetaValues} & $\zeta$-values, \code{oddIdx}, common denominator.\\
\code{Forms} & Lemmas 1--3, well-poised symmetry, $\zeta(3)$-elimination.\\
\code{DnBound} & Hanson's $d_n\le3^n$.\\
\code{Numeric} & $3^{33}g(x_0)<1$.\\
\code{Main} & root-test endgame; theorem \eqref{eq:main}.\\
\end{tabular}
\end{center}

\paragraph{Method.} The proof was assembled by an orchestration of Opus
language-model agents: the inter-module interface statements were fixed in
advance (so that the arithmetic, asymptotic, denominator and numeric parts
could be developed against stable targets), and the pieces were proved in
roughly a dozen focused rounds and merged, keeping the build green throughout.

\begin{remark}[On sharpness]
Nothing here improves on Zudilin; on the contrary, formalizing $d_n\le3^n$
rather than the prime number theorem costs the eight extra zeta values
$\zeta(27),\dots,\zeta(33)$. Once Mathlib gains the prime number theorem, the
single lemma \code{lcmUpto\_le\_three\_pow} can be replaced by
$d_n^{1/n}\to e$ and the constant $33$ throughout \code{Forms}, \code{Numeric},
\code{Main} lowered to $25$, recovering Zudilin's optimal statement with no
change to the arithmetic or asymptotic machinery.
\end{remark}

\begin{thebibliography}{9}
\bibitem{BallRivoal}
  K.~Ball, T.~Rivoal,
  \emph{Irrationalit\'e d'une infinit\'e de valeurs de la fonction z\^eta aux
  entiers impairs},
  Invent.\ Math.\ \textbf{146} (2001), 193--207.
\bibitem{Hanson1972}
  D.~Hanson,
  \emph{On the product of the primes},
  Canad.\ Math.\ Bull.\ \textbf{15} (1972), 33--37.
\bibitem{Zudilin2018}
  W.~Zudilin,
  \emph{One of the odd zeta values from $\zeta(5)$ to $\zeta(25)$ is
  irrational. By elementary means},
  SIGMA \textbf{14} (2018), 028, arXiv:1801.09895.
\end{thebibliography}

\end{document}
